\section{Veículos elétricos e cargas modernas}

\begin{frame}{Recarga de VE: por que não é uma tomada comum}
\small
\begin{columns}[T,onlytextwidth]
\column{0.50\textwidth}
\begin{itemize}
\item Potência elevada durante longos períodos.
\item Carga eletrônica com requisitos próprios de proteção.
\item Possibilidade de corrente residual CC conforme arquitetura.
\item Expansão simultânea de vários carregadores.
\item Aumento de carga pode alterar o padrão de entrada.
\item Gestão de demanda pode se tornar parte do projeto.
\end{itemize}
\column{0.48\textwidth}
\begin{caixaNorma}[title=Referências]
\begin{itemize}
\item ABNT NBR 17019:2022 — instalação fixa para alimentação/recepção de energia de VE;
\item NT.00042.EQTL Rev.04 — conexão de estações de recarga na área da Equatorial;
\item NT.00001/NT.00004 — entrada individual ou empreendimento com múltiplas UCs, conforme o caso.
\end{itemize}
\end{caixaNorma}
\end{columns}
\end{frame}

\begin{frame}{Arquitetura correta de um ponto de recarga}
\small
\centering
\begin{tikzpicture}[scale=0.78,every node/.style={font=\scriptsize}]
\tikzset{b/.style={draw=corPrim,fill=corDef!7,rounded corners,minimum width=2.25cm,minimum height=0.68cm,align=center}}
\node[b] (q) at (0,2.1) {QD};
\node[b] (p) at (3.0,2.1) {Proteção dedicada\\sobrecorrente + DR};
\node[b] (e) at (6.0,2.1) {SAVE / EVSE};
\node[b] (v) at (9.0,2.1) {Veículo};
\node[b] (pe) at (4.5,0.35) {Barramento PE / BEP\\equipotencialização};
\draw[-{Stealth},thick,corPrim] (q)--(p)--(e)--(v);
\draw[corNorma,thick] (q.south)--(pe.west);
\draw[corNorma,thick] (p.south)--(pe);
\draw[corNorma,thick] (e.south)--(pe.east);
\draw[corNorma,thick] (pe.south)--(4.5,-0.45) node[ground]{};
\end{tikzpicture}
\vspace{2mm}
\begin{caixaAt}
O PE é contínuo e integrado ao sistema de equipotencialização. Não se desenha um eletrodo de terra independente para cada equipamento. O tipo de DR deve ser compatível com a corrente residual possível, a detecção CC incorporada ao SAVE e as instruções normativas/do fabricante.
\end{caixaAt}
\end{frame}

\begin{frame}{Exemplo: SAVE monofásico de 7,4 kW}
\small
\begin{columns}[T,onlytextwidth]
\column{0.49\textwidth}
\begin{align*}
P&=\SI{7.4}{kW},\quad V=\SI{220}{V}\\
I_b&\approx\frac{7400}{220}=\SI{33.6}{A}
\end{align*}
\begin{itemize}
\item circuito dedicado;
\item cabo por ampacidade + queda + método;
\item proteção contra sobrecorrente;
\item proteção diferencial compatível;
\item PE e DPS conforme o sistema.
\end{itemize}
\column{0.49\textwidth}
\begin{caixaProjeto}[title=Impacto no projeto residencial]
No exemplo-base deste material:
\[
33{,}96+7{,}40=\SI{41.36}{kW}.
\]
Logo, a inclusão do SAVE não é apenas ``mais um circuito'': no exemplo, a faixa da Tabela 1 da NT.00001 passa de 24,1--38 kW (63 A) para 38,1--48 kW (80 A). QD, alimentador, balanceamento e solicitação de aumento de carga devem ser revistos.
\end{caixaProjeto}
\end{columns}
\end{frame}

\begin{frame}{Gestão dinâmica de recarga}
\small
\centering
\begin{tikzpicture}[scale=0.80,every node/.style={font=\scriptsize}]
\tikzset{b/.style={draw=corPrim,fill=corDef!7,rounded corners,minimum width=2.25cm,minimum height=0.65cm,align=center}}
\node[b] (med) at (0,2.2) {Medição da demanda\\da instalação};
\node[b] (ctrl) at (3.2,2.2) {Controlador\\de carga};
\node[b] (ev1) at (6.4,3.2) {SAVE 1};
\node[b] (ev2) at (6.4,2.2) {SAVE 2};
\node[b] (ev3) at (6.4,1.2) {SAVE 3};
\node[b] (pred) at (0,0.5) {Demais cargas\\da edificação};
\draw[-{Stealth},thick,corPrim] (med)--(ctrl);
\foreach \x in {ev1,ev2,ev3}{\draw[-{Stealth},thick,corAccent] (ctrl)--(\x);}
\draw[-{Stealth},thick,corPrim] (med)--(pred);
\node[align=center] at (3.2,0.45) {limite disponível\\+ prioridades};
\end{tikzpicture}
\vspace{1mm}
\begin{caixaObs}
Gestão de carga reduz simultaneidade dos carregadores, mas não elimina as verificações de capacidade térmica, proteção, qualidade de energia e requisitos da distribuidora.
\end{caixaObs}
\end{frame}

\begin{frame}{Geração distribuída: interface com a instalação}
\small
\begin{columns}[T,onlytextwidth]
\column{0.49\textwidth}
\begin{itemize}
\item Corrente pode fluir em sentido diferente do projeto tradicional.
\item Proteções e seccionamento existem nos lados CC e CA.
\item DPS e equipotencialização precisam ser compatíveis com o SPDA.
\item Verificar capacidade de barramentos e condutores no ponto de conexão.
\item O estudo interno deve fechar com a solicitação de conexão à distribuidora.
\end{itemize}
\column{0.49\textwidth}
\begin{caixaNorma}[title=Equatorial]
A conexão de micro e minigeração distribuída é tratada pela \textbf{NT.00020.EQTL Rev.06}. Para unidades do Grupo B, usar também os formulários e anexos atuais disponibilizados pela distribuidora.
\end{caixaNorma}
\begin{caixaAt}
Geração distribuída não autoriza reduzir proteção, seção de condutor ou capacidade de barramento apenas porque parte da energia é produzida localmente.
\end{caixaAt}
\end{columns}
\end{frame}

\begin{frame}{UPS, geradores e transferência}
\small
\begin{columns}[T,onlytextwidth]
\column{0.49\textwidth}
\begin{itemize}
\item Definir cargas essenciais e não essenciais.
\item Analisar curto-circuito em modo rede e modo fonte alternativa.
\item A proteção que atua bem na rede pode não atuar igual com fonte limitada eletronicamente.
\item Verificar neutro, aterramento e chaveamento conforme a topologia.
\end{itemize}
\column{0.49\textwidth}
\centering
\begin{tikzpicture}[scale=0.78,every node/.style={font=\scriptsize}]
\tikzset{b/.style={draw=corPrim,fill=corDef!7,rounded corners,minimum width=2.1cm,minimum height=0.62cm,align=center}}
\node[b] (r) at (0,2.6) {Rede};
\node[b] (g) at (0,0.8) {Gerador};
\node[b] (ats) at (3.2,1.7) {ATS};
\node[b] (ups) at (6.0,1.7) {UPS};
\node[b] (crit) at (8.8,1.7) {Cargas críticas};
\draw[-{Stealth},thick,corPrim] (r)--(ats);\draw[-{Stealth},thick,corAccent] (g)--(ats);\draw[-{Stealth},thick,corPrim] (ats)--(ups)--(crit);
\end{tikzpicture}
\end{columns}
\end{frame}

%=====================================================================
